\documentclass[a4paper,10pt]{article}

\usepackage[utf8]{inputenc}
\usepackage{t1enc}
\usepackage[spanish]{babel}
\usepackage[pdftex,usenames,dvipsnames]{color}
\usepackage[pdftex]{graphicx}
\usepackage{amsmath}
\usepackage{amsfonts}
\usepackage{amssymb}
\usepackage{listings}
\lstset{language=C}
\lstset{showstringspaces=false}
\lstset{basicstyle=\ttfamily,}

\begin{document}

\begin{titlepage}
        \thispagestyle{empty}
        \begin{center}
                \includegraphics{./images/itba.jpg}
                \vfill
                \Huge{Arquitectura de las Computadoras}\\
                \vspace{1cm}
                \huge{Trabajo Práctico Especial}\\
        \end{center}
        \vspace{2cm}
        \large{
                \begin{tabular}{lcrc}
                        \textbf{Alvaro Crespo} & & 47283 & \ \ \texttt{acrespo@alu.itba.edu.ar}\\
                        \textbf{Juan Pablo Civile} & & 50753 & \ \ \texttt{jcivile@alu.itba.edu.ar}\\
                        \textbf{Dario Susnisky} & & 48083 & \ \ \texttt{dsusnisk@alu.itba.edu.ar}\\
                        \\ 
                \end{tabular}
        }
        \vfill
        \flushright{1 de Junio del 2011}
\end{titlepage}

\setcounter{page}{1}

\tableofcontents
\newpage

\section{Resumen}

\section{Diseño del Kernel}
    Las funcionalidades del kernel fueron desarrolladas incrementalmente. Inicialmente, optamos por implementar las system calls read y write de una manera muy sencilla.
    Una vez que teniamos esta version desarrollada, comenzamos a pensar como interactuan estas con las aplicaciones. 
    Al estar usando read y write como las define Linux, decidimos investigar el I/O de teclado y video en las terminales de sistemas Unix-compatible.
    Dado que el objetivo no del desarrollo no es diseñar un sistema operativo nuevo, optamos por ser lo mas compatibles con Linux posible.

    Investigando aprendimos que Linux maneja esta interaccion de una manera que emula las viejas terminales de mainframes. 
    Esto significa, que hay solo 3 operaciones posibles: read, write e ioctl.
    ioctl es una system call que permite comunicarse con los drivers detras de un dado archivo de maneras especificas a cada driver.
    Esta es la llamada que permite manipular el IO para obtener los distintos comportamientos observados en una consola en Linux.
    Tambien encontramos que las terminales aceptan las llamadas escape sequences, que permiten realizar operaciones especiales sobre la pantalla, como manipular el color y el cursor.

    En fin, esto derivo en la creacion de una emulacion de una tty a nivel kernel.
    Si bien la emulacion es limitada, es en parte compatible con la mayoria de las shells modernas de Linux.
    Para esto se crearon los drivers de video y teclado, con los cuales se comunica desde las aplicaciones usando read, write e ioctl.
    Dado que no existen mas drivers, read y write se comunican directamente con los drivers.

    Esta emulacion nos permitio gran flexibilidad y aporto un nivel de funcionalidad muy alto para el desarrollo de aplicaciones.
    Por ejemplo, se presto con facilidad a la implementacion de multiples tty.

    \subsection{Interrupciones}

    \subsection{Terminales}
        \subsubsection{Entrada}
            La entrada se maneja mediante un buffer, que es actualizado cuando el PIC dispara una interrupcion de parte del controlador de teclado.
            Cuando esta interrupcion se dispara, se interpreta los codigos enviados por el teclado y se agregan al buffer.
            La interpretacion de estos codigos, difiere segun el estado de ciertas teclas, como serlo Alt, Ctrl, Shift y Bloc Mayus.
            Luego, este buffer puede ser consultado utilizando la llamada read sobre la entrada standard.

            Como fue mencionado anteriormente, las terminales soportan varios modos de entrada.
            Por simplicidad, y por que proveen suficiente funcionalidad, se implementaron 4 modos en total.
            Que modo se utiliza se puede manipular con la llamada ioctl.

            El modo esta determinado por 2 flags: canon y echo.
            Canon activa y desactiva el modo Canonico de la terminal.
            Mientra que echo activa y desactiva si los caracteres leidos son impresos a pantalla por la terminal.

            El modo canonico determina un gran factor del comportamiento de llamar a read sobre la entrada standard.
            Si este esta activado, read funciona como buffereado por linea, o sea, lee de a lineas.
            Esto significa que en una dada llamada, read va a retornar como mucho una linea, y no va a retornar, hasta que tenga por lo menos una linea en su buffer.
            Si bien las terminales de linux permiten manipular que determina un fin de linea, nosotros tomamos \n como fin de linea.
            Si el modo canonico no esta activado, read se comporta como su definicion dice, lee n bytes, sin importar si hay lineas o no.

            Ademas de soportar estos 4 modos, el driver interpreta teclas especiales, como lo son las flechas.
            Algunas de estas teclas, como ser Bloc Mayus, Num Lock y Scroll Lock son manejadas internamente por el driver.
            Otras como las flechas, Inicio, Fin y las teclas de Funcion, se guardan como parte de la entrada previa conversion a escape sequences.

            Por limitaciones del sistema, al no tener multi tarea, el buffer de teclado es uno solo, compartido por todas las terminales.

        \subsubsection{Salida}
            La comunicacion entre el driver de salida de una terminal y la controladora de video se hace mediante el sector de video mapeado a memoria.
            Este permite suficiente flexibilidad para nuestros objetivos, y a su vez es muy simple de utilizar.
            La unica excepcion es la manipulacion del cursor grafico, que se hace utilizando in y out, ya que no encontramos si la posicion del mismo esta mapeada en alguna parte de memoria.

            El driver permite manipular todos sus aspectos mediante llamadas a write.
            Este acepta escape sequences que dejan realizar un gran abanico de operaciones sobre la salida.
            El listado de esta se puede encontrar en el anexo <<Completar aca>>.
            
            En este nivel se maneja un buffer de salida donde se guarda una copia del estado actual de la pantalla.
            Esto es asi por 2 motivos.
            Primero, para evitar lecturas a la memoria de video cuando es necesario hacer scroll de pantalla.
            Segundo, es un elemento crucial a la hora de implementar multiples terminales.
            Cada terminal tiene su propio estado en el driver de salida, incluyendo su propio buffer, a diferencia del driver de entrada.
            Cuando es pedido, el driver cambia que estado usa, efectivamente cambiando de terminal.

\section{Libreria standard}

\section{Shell}
    Contamos con una shell por cada terminal.
    Esta utiliza al maximo las capacidades de la terminal, comunicandose siempre por la libreria standard.
    La shell manipula el modo de la entrada de la terminal para permitirse manipular la entrada con gran precision.
    Se desactiva tanto el echo como el modo canonico, o sea, lo ingresado llega solo a la shell y sin esperar a que se complete la linea.
    Esto le permite al ususario editar su comando a medida que lo escribe, auto completar y poder moverse por el historial de comandos.

    El comando sudoku demuestra la capacidad de la terminal, haciendo uso de los colores y la posibilidad de mover el cursor a placer.

    Se menciono antes que por la limitacion de ser mono tarea, hay un solo buffer de teclado.
    Esto se nota tambien a nivel shell ya que solo se permite cambiar entre terminales cuando se esta ingresando un comando.

\end{document}
